\chapter{Tổ hợp}

\section{Hoán vị}
	\subsection{Giai thừa}
		\import{factorial.tex}
		\kactlimport{IntPerm.h}

	\subsection{Chu trình}
		Gọi $g_S(n)$ là số hoán vị bậc $n$ mà độ dài mọi chu trình đều thuộc tập $S$. Khi đó:
		$$\sum_{n=0} ^\infty g_S(n) \frac{x^n}{n!} = \exp\left(\sum_{n\in S} \frac{x^n} {n} \right)$$

	\subsection{Hoán vị mất thứ tự (Derangements)}
		Số hoán vị của $n$ phần tử sao cho không phần tử nào ở đúng vị trí ban đầu:
		\[ \mkern-2mu D(n) = (n-1)(D(n-1)+D(n-2)) = n D(n-1)+(-1)^n = \left\lfloor\frac{n!}{e}\right\rceil \]

	\subsection{Bổ đề Burnside}
		Cho nhóm đối xứng $G$ tác động lên tập $X$, số lớp tương đương (cấu hình phân biệt) theo phép đối xứng bằng:
		 \[ {\frac {1}{|G|}}\sum _{{g\in G}}|X^{g}|, \]
		 trong đó $X^{g}$ là tập các phần tử bất biến dưới phép biến đổi $g$ ($g.x = x$).

		 Nếu xét các cấu hình độ dài $n$ bất biến qua phép quay ($G = \mathbb Z_n$):
		 \[ g(n) = \frac 1 n \sum_{k=0}^{n-1}{f(\gcd(n, k))} = \frac 1 n \sum_{k|n}{f(k)\phi(n/k)}. \]

\section{Phân hoạch và tập con}
	\subsection{Hàm phân hoạch (Partition function)}
		Số cách phân tích $n$ thành tổng các số nguyên dương không xét thứ tự:
		\[ p(0) = 1,\ p(n) = \sum_{k \in \mathbb Z \setminus \{0\}}{(-1)^{k+1} p(n - k(3k-1) / 2)} \]
		\[ p(n) \sim 0.145 / n \cdot \exp(2.56 \sqrt{n}) \]

		\begin{center}
		\begin{tabular}{c|c@{\ }c@{\ }c@{\ }c@{\ }c@{\ }c@{\ }c@{\ }c@{\ }c@{\ }c@{\ }c@{\ }c@{\ }c}
			$n$    & 0 & 1 & 2 & 3 & 4 & 5 & 6  & 7  & 8  & 9  & 20  & 50  & 100 \\ \hline
			$p(n)$ & 1 & 1 & 2 & 3 & 5 & 7 & 11 & 15 & 22 & 30 & 627 & $\mathtt{\sim}$2e5 & $\mathtt{\sim}$2e8 \\
		\end{tabular}
		\end{center}

	\subsection{Định lý Lucas}
		Cho $n,m$ là các số nguyên không âm và $p$ là số nguyên tố. Viết $n=n_kp^k+...+n_1p+n_0$ và $m=m_kp^k+...+m_1p+m_0$ trong hệ cơ số $p$. Khi đó:
		\[ \binom{n}{m} \equiv \prod_{i=0}^k\binom{n_i}{m_i} \pmod{p} \]

	\subsection{Hệ số nhị thức \& Đa thức}
		\kactlimport{multinomial.h}

\section{Các dãy số đặc biệt}
	\subsection{Số Bernoulli}
		Hàm sinh mũ của số Bernoulli: $B(t)=\frac{t}{e^t-1}$ (tính được bằng FFT).
		$B[0,\ldots] = [1, -\frac{1}{2}, \frac{1}{6}, 0, -\frac{1}{30}, 0, \frac{1}{42}, \ldots]$

		Tổng lũy thừa:
		\small
		\[ \sum_{i=1}^n n^m = \frac{1}{m+1} \sum_{k=0}^m \binom{m+1}{k} B_k \cdot (n+1)^{m+1-k} \]
		\normalsize

	\subsection{Số Stirling loại một}
		Số hoán vị của $n$ phần tử có đúng $k$ chu trình rời nhau:
		\begin{align*}
			&c(n,k) = c(n-1,k-1) + (n-1) c(n-1,k),\ c(0,0) = 1 \\
			&\textstyle \sum_{k=0}^n c(n,k)x^k = x(x+1) \dots (x+n-1)
		\end{align*}

	\subsection{Số Euler (Eulerian numbers)}
		Số hoán vị $\pi \in S_n$ có đúng $k$ vị trí tăng $\pi(j) < \pi(j+1)$:
		$$E(n,k) = (n-k)E(n-1,k-1) + (k+1)E(n-1,k)$$
		$$E(n,0) = E(n,n-1) = 1$$

	\subsection{Số Stirling loại hai}
		Số cách phân hoạch $n$ phần tử phân biệt vào đúng $k$ nhóm không rỗng:
		$$S(n,k) = S(n-1,k-1) + k S(n-1,k)$$
		$$S(n,1) = S(n,n) = 1$$
		$$S(n,k) = \frac{1}{k!}\sum_{j=0}^k (-1)^{k-j}\binom{k}{j}j^n$$

	\subsection{Số Bell}
		Tổng số cách phân hoạch $n$ phần tử phân biệt vào các nhóm. $B(n) =$
		$1, 1, 2, 5, 15, 52, 203, 877, 4140, 21147, \dots$.

	\subsection{Đếm cây có nhãn (Cayley's formula / Prüfer)}
		Số cây trên $n$ đỉnh: $n^{n-2}$ \\
		Số cách nối $k$ cây kích thước $n_i$ thành 1 cây: $n_1n_2\cdots n_k n^{k-2}$ \\
		Số cây với bậc các đỉnh là $d_i$: $(n-2)! / ((d_1-1)! \cdots (d_n-1)!)$

	\subsection{Số Catalan}
		\[ C_n=\frac{1}{n+1}\binom{2n}{n}= \binom{2n}{n}-\binom{2n}{n+1} = \frac{(2n)!}{(n+1)!n!} \]
		\[ C_0=1,\ C_{n+1} = \frac{2(2n+1)}{n+2}C_n,\ C_{n+1}=\sum C_iC_{n-i} \]
		${C_n = 1, 1, 2, 5, 14, 42, 132, 429, 1430, 4862, 16796, 58786, \dots}$
		\begin{itemize}[noitemsep]
			\item Số đường đi đơn điệu dưới đường chéo trên lưới $n\times n$.
			\item Dãy ngoặc đúng độ dài $2n$ gồm $n$ cặp ngoặc.
			\item Số cây nhị phân có $n+1$ lá ($0$ hoặc $2$ con).
			\item Số cây có gốc có thứ tự gồm $n+1$ đỉnh.
			\item Số cách tam giác hóa đa giác lồi $n+2$ đỉnh bằng các đường chéo không cắt nhau.
			\item Số hoán vị của $[n]$ không chứa dãy con tăng độ dài 3.
		\end{itemize}
